\documentclass[11pt, a4paper]{article}
\usepackage[utf8]{inputenc}
\usepackage{geometry}
\geometry{margin=1in}
\usepackage{graphicx}
\usepackage{amsmath}
\usepackage{booktabs}
\usepackage{longtable}
\usepackage{xcolor}
\usepackage{enumitem}
\usepackage{fancyhdr}
\usepackage{titlesec}
\usepackage{hyperref}
\hypersetup{colorlinks=true, linkcolor=blue!70!black, urlcolor=blue!70!black}

% Header/Footer
\pagestyle{fancy}
\fancyhf{}
\fancyhead[L]{\textsc{TWP Logbook — Assignment 2A}}
\fancyhead[R]{\textsc{CSP529}}
\fancyfoot[C]{\thepage}
\renewcommand{\headrulewidth}{0.4pt}

\titleformat{\section}{\large\bfseries}{\thesection.}{0.5em}{}[\titlerule]
\titleformat{\subsection}{\normalsize\bfseries}{\thesubsection.}{0.5em}{}

\title{
    \vspace{-1cm}
    \textbf{Writing Logbook --- Assignment 2A} \\[0.3em]
    \large CSP529 Technical Writing \& Publishing \\[0.2em]
    \large \textit{From \LaTeX\ Elements to a Professional Technical Document}
}
\author{
    \textbf{MT26MCS024} --- Gaurav Acharya \quad
    \textbf{MT26MCS032} --- Sandeep
}
\date{September 2026}

\begin{document}
\maketitle
\thispagestyle{fancy}

% ============================================================
\section{Assignment Overview}
% ============================================================

\begin{tabular}{@{}p{4.5cm} p{10cm}@{}}
    \toprule
    \textbf{Item} & \textbf{Detail} \\
    \midrule
    Practical Title & Practical 2A --- From \LaTeX\ Elements to a Professional Technical Document \\
    Topic & Restructuring the OS Process document into a proper academic paper \\
    Group Members & Gaurav Acharya (MT26MCS024), Sandeep (MT26MCS032) \\
    Builds Upon & Assignment 02 (Stage 2: Representations of an OS Process) \\
    Date of Completion & September 2026 \\
    \bottomrule
\end{tabular}

% ============================================================
\section{Objectives \& Learning Outcomes}
% ============================================================

This assignment extended Assignment~02 by transforming a collection of isolated \LaTeX\ elements into a cohesive, professionally structured technical document. The learning outcomes included:

\begin{itemize}[nosep]
    \item Restructuring content into standard academic sections: Title $\rightarrow$ Abstract $\rightarrow$ Introduction $\rightarrow$ Technical Background $\rightarrow$ Analysis/Working $\rightarrow$ Discussion $\rightarrow$ Conclusion $\rightarrow$ References.
    \item Mastering floats: captions, labels, and cross-references for figures and tables.
    \item Making equations part of the argument, not decorative additions.
    \item Understanding that algorithms describe \emph{procedure} while prose communicates \emph{purpose, rationale, and interpretation}.
    \item Learning \LaTeX\ cross-referencing (\texttt{\textbackslash ref\{\}}), citation mechanics (\texttt{\textbackslash cite\{\}}), and bibliography management.
\end{itemize}

% ============================================================
\section{Work Performed --- Stage by Stage}
% ============================================================

\subsection{Stage 1: Inspecting Previous Work (Practical 2)}
We critically examined our Assignment~02 document and identified:
\begin{itemize}[nosep]
    \item \LaTeX\ constructs used: \texttt{figure}, \texttt{table}, \texttt{equation}, \texttt{algorithm}, \texttt{itemize}, \texttt{enumerate}.
    \item Elements that were difficult: float placement, equation formulation.
    \item Elements inserted only because required: some lists felt mechanical rather than communicative.
    \item Elements that genuinely helped: the equation and algorithm provided clarity that prose alone could not.
    \item \textbf{Critical finding:} Most elements were \emph{not} referred to from the body text --- they simply ``appeared'' without introduction or interpretation.
\end{itemize}

\noindent\textbf{Key Rule Discovered:} \textit{Nothing should appear in a technical document without being introduced, referred to, or interpreted in the text.}

\subsection{Stage 2: Building Proper Document Structure}
We restructured the entire document into:
\begin{enumerate}[nosep]
    \item \textbf{Abstract:} A concise overview of the document's scope --- OS boot process, hardware-software interfacing, kernel resource management.
    \item \textbf{Introduction:} Context and motivation, with an inline citation~\cite{silberschatz}.
    \item \textbf{Technical Background:} Component comparison table (now with textual reference: ``The comparison presented in Table~1 indicates\ldots'') and the state-transition equation (now with interpretation).
    \item \textbf{Analysis/Working:} Chronological boot sequence (numbered list), algorithmic pseudocode, and the architectural figure (now with: ``The computational steps are illustrated in Figure~1'').
    \item \textbf{Discussion:} Plain-language explanation of the algorithm, plus OS core responsibilities.
    \item \textbf{Conclusion:} Summary of findings.
    \item \textbf{References:} Three genuine academic sources.
\end{enumerate}

\textbf{Justification:} Technical Background establishes definitions and models \emph{before} they are used in Analysis. Discussion interprets the formal constructs in accessible language.

\subsection{Stage 3: Mastering Figures and Tables}
We ensured every float was properly integrated:
\begin{itemize}[nosep]
    \item Figure with caption (``Operating System Architecture Layers''), label (\texttt{fig:os\_process}), and cross-reference from text.
    \item Table with caption (``System Component and Execution Phase''), label (\texttt{tab:os\_components}), and cross-reference from text.
\end{itemize}

\noindent\textbf{Discoveries:}
\begin{itemize}[nosep]
    \item Figures and tables are numbered automatically by \LaTeX.
    \item Inserting a new figure before an existing one automatically renumbers all subsequent figures.
    \item \texttt{\textbackslash ref\{\}} ensures references update automatically; manual numbering would break on any restructuring.
    \item Captions are essential because they allow figures/tables to be understood independently of surrounding text.
\end{itemize}

\subsection{Stage 4: Mathematics as Part of the Argument}
We included two meaningful equations:
\begin{enumerate}[nosep]
    \item \textbf{Equation~1 --- OS State Transition:} $S_{t+1} = k\left(S_t, \sum I_i + \sum C_j\right) - \Delta_{overhead}$
    
    Numbered, labelled, cross-referenced, and interpreted: ``This equation demonstrates that the OS continuously evaluates the current state alongside hardware interrupts and system calls\ldots''
    
    \item \textbf{Equation~2 --- Turnaround Time:} $T_{turnaround} = T_{completion} - T_{arrival}$
    
    Added as a second equation to analyse CPU scheduling efficiency, also numbered, labelled, and interpreted.
\end{enumerate}

\noindent\textbf{Answer to Stage 4 question:} The equations communicate \emph{quantitative relationships} and \emph{variable dependencies} that the paragraph alone cannot express with the same precision and conciseness.

\subsection{Stage 5: Algorithm vs.\ Explanation}
We maintained the OS Boot Sequence pseudocode and added explanatory prose in the Discussion section that translates the algorithmic steps into plain language.

\noindent\textbf{Discovery:} Pseudocode describes \emph{what happens step by step} (procedure), while prose communicates \emph{why} each step matters, what assumptions underlie it, and how to interpret the outcomes (purpose and rationale).

\subsection{Stage 6: Cross-Referencing Challenge}
We added cross-references to: 2 sections, 2 figures/tables, 1 equation, and 1 algorithm. After inserting a new section and figure in the middle of the document and recompiling:
\begin{itemize}[nosep]
    \item \textbf{What changed automatically:} All section numbers, figure numbers, table numbers, and equation numbers updated throughout the entire document.
    \item \textbf{Why logical references matter:} In a professional document with dozens of figures, manually updating ``Figure~3'' to ``Figure~4'' everywhere would be error-prone and unsustainable. \texttt{\textbackslash ref\{\}} ensures consistency by design.
\end{itemize}

\subsection{Stage 7: References}
We included three genuine references connected to our technical content:
\begin{enumerate}[nosep]
    \item Silberschatz, Galvin, Gagne --- \textit{Operating System Concepts}, 10th ed., 2018.
    \item Tanenbaum, Bos --- \textit{Modern Operating Systems}, 4th ed., 2014.
    \item Love --- \textit{Linux Kernel Development}, 3rd ed., 2010.
\end{enumerate}

\noindent\textbf{Skills learned:} \texttt{\textbackslash cite\{\}}, \texttt{thebibliography} environment, citation-reference distinction.

\subsection{Stage 8: \LaTeX\ Debugging Challenge}

\begin{tabular}{@{}p{4cm} p{5cm} p{5cm}@{}}
    \toprule
    \textbf{Error/Problem} & \textbf{What Caused It} & \textbf{How It Was Corrected} \\
    \midrule
    Figure not appearing & Missing \texttt{float} package for \texttt{[H]} placement & Added \texttt{\textbackslash usepackage\{float\}} \\
    Undefined reference ``??'' & \texttt{\textbackslash label} placed before \texttt{\textbackslash caption} & Moved \texttt{\textbackslash label} after \texttt{\textbackslash caption} \\
    Missing package error & Used \texttt{array} column types without importing & Added \texttt{\textbackslash usepackage\{array\}} \\
    \bottomrule
\end{tabular}

% ============================================================
\section{Comparison: Assignment 02 vs.\ Assignment 2A}
% ============================================================

\begin{tabular}{@{}p{4cm} p{5cm} p{5cm}@{}}
    \toprule
    \textbf{Aspect} & \textbf{Assignment 02 (Before)} & \textbf{Assignment 2A (After)} \\
    \midrule
    Structure & Flat list of representations & Full academic structure (Abstract to References) \\
    Floats & Inserted without context & Introduced, referenced, and interpreted in text \\
    Equations & Single equation, no interpretation & Two equations, both numbered, labelled, interpreted \\
    References & None & Three genuine academic sources with \texttt{\textbackslash cite\{\}} \\
    Cross-references & None & Sections, figures, tables, equations all cross-referenced \\
    \bottomrule
\end{tabular}

% ============================================================
\section{Key Takeaways}
% ============================================================

\begin{enumerate}[nosep]
    \item The transition from ``using \LaTeX\ elements'' to ``writing a technical document'' requires every element to serve a communicative purpose within the document's argument.
    \item Cross-referencing is not a convenience --- it is a requirement for any document that may be revised or extended.
    \item The \texttt{Abstract $\rightarrow$ Introduction $\rightarrow$ Background $\rightarrow$ Analysis $\rightarrow$ Discussion $\rightarrow$ Conclusion} structure is not arbitrary; it mirrors the reader's cognitive needs.
    \item Debugging \LaTeX\ compilation errors (undefined references, missing packages, float placement) is a core practical skill that improves with experience.
\end{enumerate}

% ============================================================
% Dummy bibliography for \cite references in this logbook
\begin{thebibliography}{9}
\bibitem{silberschatz} A.\ Silberschatz, P.\ B.\ Galvin, G.\ Gagne, \textit{Operating System Concepts}, 10th ed., Wiley, 2018.
\end{thebibliography}

\end{document}
